\documentclass[11pt]{article}
\usepackage[a4paper,margin=1in]{geometry}
\usepackage[T1]{fontenc}
\usepackage{lmodern,amsmath,amssymb,amsthm,booktabs,microtype}
\usepackage[hidelinks]{hyperref}
\hypersetup{pdftitle={Counting Defects: Semilinear Reachability and Resource Bounds for Langton's Ant},pdfauthor={Hao Ke}}
\usepackage{enumitem}
\setlist{nosep}
\setlength{\emergencystretch}{2em}
\newtheorem{theorem}{Theorem}[section]
\newtheorem{lemma}[theorem]{Lemma}
\newtheorem{corollary}[theorem]{Corollary}
\newtheorem{proposition}[theorem]{Proposition}
\theoremstyle{remark}\newtheorem{remark}[theorem]{Remark}
\newcommand{\Z}{\mathbb Z}\newcommand{\N}{\mathbb N}
\newcommand{\Rr}{\mathrm R}\newcommand{\Ll}{\mathrm L}
\newcommand{\dom}{\operatorname{dom}}\newcommand{\ran}{\operatorname{ran}}
\title{Counting Defects: Semilinear Reachability\\and Resource Bounds for Langton's Ant}
\author{Hao Ke}
\date{}
\begin{document}
\maketitle
\begin{abstract}
Fix a number of initially black cells and a target number of departures
from the standard Langton highway. Which placements of those cells make the
ant reach that target? We count a departure whenever ten consecutive turns
form a word absent from the standard 104-turn cycle. For each supplied pair
of integers $m,N$, we prove that the placements of exactly $m$ black cells
producing at least $N$ such departures form an effectively computable
semilinear set: a finite union of integer translates of finitely generated
additive monoids. Initial coordinates are unrestricted. The construction also
describes the exact time, pose and entire black field just after the $N$th
departure. The key reduction represents the intervening defect-free motion by a
finite collection of words with variable repetition counts. Consecutive visits
to each cell enforce the ant's read-and-flip rule throughout these compressed
traces. This yields decidable synthesis under integer-linear coordinate
constraints and piecewise rational-affine resource formulas for supplied
word families. Positive instances have polynomial-bit witnesses when $m,N$
are written in unary. A restricted quantifier-free encoding makes the
spatial read/write tests polynomial in the supplied description size.
Together these results turn finite departure targets into exact arithmetic
questions about complete ant states.
A fixed six-cell initial shape supplies an unbounded resource ray:
translation by a parameter $k$ produces exactly $6702+47k$ forward defects
before permanent highway entrance at $18743+220k$. Consequently every finite
defect target can be reached with at most six initially black cells.
\end{abstract}

\section{Problem and main result}
Suppose we may place $m$ black cells anywhere on an otherwise white plane.
We want the ant to depart from the local turn patterns of its standard highway
at least $N$ times. The cell coordinates can be arbitrarily large, and the
$N$th departure may occur after an arbitrarily long flight. This is therefore
a synthesis problem over unbounded coordinates and time. We show that each
fixed target $m,N$ has an exact description in integer-linear arithmetic.

Here \emph{Presburger arithmetic} means first-order arithmetic with integer
addition, order and fixed integer constants, but without multiplication of
two variables. Its definable sets are exactly the semilinear sets: finite
unions of sets $a+\mathbb N v_1+\cdots+\mathbb N v_r$ with integer vectors
$a,v_i$. Such a description allows one to test whether any initial placement
meets additional coordinate constraints. Our construction also records the
complete state when the target is first reached.
A state is $(B,p,h)$, where $B\subset\Z^2$ is the \emph{complete} finite set of
black cells, $p\in\Z^2$ is the head position, and $h\in\{N,E,S,W\}$.
Encode these headings by the integers $0,1,2,3$, respectively. A white read turns
right, a black read turns left, the current cell flips, and the ant then moves
one unit. Time $t$ denotes the state immediately before update $t$.

Record a right turn by R and a left turn by L. Let $W$ be the standard
period-104 turn word printed in Appendix~\ref{app:words},
and let $\mathcal L_{10}$ be its cyclic ten-letter factors. A \emph{contained
defect} of a finite forward word is a length-ten consecutive subword outside
$\mathcal L_{10}$. For example, a segment of the repeated highway word has
no defects, even across the boundary between two copies. A contained defect
is indexed by its final update $t\ge9$; its post-event state is at time $t+1$.
The ten-factor
attraction theorem states that a finite-support orbit is eventually standard
if and only if it has finitely many such defects \cite{project}.

Normalize the head to the origin facing north. For fixed $m\ge0$ and $N\ge1$,
let $E_{m,N}$ contain the ordered, pairwise-distinct tuples
$z=(x_1,y_1,\ldots,x_m,y_m)$ for which the actual ant with exactly these black
cells has at least $N$ contained defects. Coordinates have no range restriction.
For $m=0$, the tuple is empty and the relation asks this question of the
all-white initial plane.

A finite-fibre field family is a definable set of pairs $(z,q)$ such that,
for each parameter tuple $z$, only finitely many lattice cells $q$ belong to
the field. Describing this membership relation specifies every endpoint cell,
including unvisited initial inventory.
For later resource questions, write
\[
 \mu(N)=\min\{|B|:\text{the complete finite input }B
              \text{ produces at least }N\text{ contained defects}\}.
\]
The initial pose is normalized as above. This quantity asks for the least
black inventory needed to reach a finite target.

\begin{theorem}[All-coordinate fixed-budget synthesis]\label{thm:main}
There is an effective construction, uniform in the supplied integers $m,N$,
of a Presburger formula defining $E_{m,N}$. The exact time and pose immediately
after the $N$th defect have an effectively Presburger-definable graph. The
complete black field at that endpoint is an effective finite-fibre Presburger
family over the initial coordinates.
\end{theorem}

In particular, existence, nonexistence, and intersection with a prescribed
Presburger coordinate domain are decidable. A finite union handles ``at most
$m$ black cells.'' This is not one fixed formula in which arbitrary $m,N$ are
additional free variables. The finite union can be large; its effective construction establishes
decidability, while the complexity results below concern compressed positive
witnesses and a restricted encoding.

The quantity $N$ is a target for a finite prefix: the ant may have further
defects after the reported endpoint. The formula describes inputs reaching
the $N$th defect and says nothing about a bound on the total future count.

The central difficulty is physical execution: the ant carries an unbounded
mutable spatial store, so a short word description must retain every earlier
read and every initially black cell. Consecutive same-cell relations solve
this problem exactly. A finite ten-factor cover then yields the all-coordinate
construction, and unary time-fibre counting yields its resource refinement.
Sections~\ref{sec:power}--\ref{sec:cover} build the arithmetic description
and prove that its finite family of word forms covers every $N$th-event
prefix. Sections~\ref{sec:supply} and~\ref{sec:small} derive resource spectra
and small witnesses; Section~\ref{sec:qf} gives the restricted quantifier-free
encoding. A \emph{word schema} means a fixed sequence of literal words and
repeated words whose nonnegative repetition counts are parameters.

Symbolic feasibility for ant words has an established history~\cite{gajardo}.
The equivalence of Presburger descriptions and semilinear sets is classical
\cite{ginsburg}. Flat-system verification already provides run schemas,
small-witness arguments and quantifier-free counting encodings~\cite{flat,demri}.
Our reduction supplies the additional physical semantics of the ant's complete
mutable black field. The companion factor paper~\cite{project} supplies the
bounded-fold exclusion and complete finite cover. Together these yield exact
all-coordinate synthesis, ant-specific supply spectra and unary-budget NP
membership. The restricted backend encodes spatial read/write parity as well
as factor counts in polynomial size. Unary counting and piecewise-affine
Presburger functions are established tools~\cite{counting,cluckers}; the
explicit fold--consumer spectrum is an application with sharp ant-specific
constants.

The construction in Section~\ref{sec:six} makes the fixed-cost branch of the
resource spectrum explicit. Its complete-field recurrence carries the ant
through arbitrarily many return words, then into the standard highway. The
result reduces the unconstrained minimum-cost search to masses zero through
six, while the synthesis theorem handles coordinates and prescribed constraints.

\section{Consecutive reads without a parity-counting oracle}
Fix a parameter tuple $\theta$ in a Presburger domain $D$. Suppose a prescribed
kinematic trace is represented by a Presburger function $V_\theta$ from a
contiguous interval of update times to physical cells. Its domain is either
$\{0,\ldots,T(\theta)-1\}$ or $\N$. Let $R_\theta,L_\theta$ partition its times
according to the requested letters. The movements and headings are assumed to
satisfy the literal turn rule; Section~\ref{sec:power} constructs such traces.

Define
\[
 C_\theta=\{(s,t):s<t,\ V_\theta(s)=V_\theta(t)\}.
\]
Relational composition uses the same parameter tuple throughout.
\begin{lemma}[Successor subtraction]\label{lem:next}
The relation
\begin{equation}\label{eq:next}
 \operatorname{Next}_\theta=C_\theta\setminus(C_\theta\circ C_\theta)
\end{equation}
consists exactly of consecutive visits to the same physical cell.
\end{lemma}
\begin{proof}
A pair in $C_\theta$ fails to be consecutive precisely when a read time $u$
strictly between its endpoints visits their common cell. This is exactly
membership in $C_\theta\circ C_\theta$. No transitive closure is computed.
\end{proof}

Put $\mathcal T_\theta=\dom V_\theta$ and
\[
 \operatorname{First}_\theta=\mathcal T_\theta\setminus\ran C_\theta.
\]
For a complete initial black predicate $B(\theta,x,y)$, write
$\operatorname{Hits}_\theta=\{t:B(\theta,V_\theta(t))\}$.

\begin{theorem}[Exact realization]\label{thm:semantics}
The parameter assignments realizing the prescribed trace are exactly those
satisfying
\begin{align}
\operatorname{Next}_\theta\cap
 ((R_\theta\times R_\theta)\cup(L_\theta\times L_\theta))&=\varnothing,
                                                        \label{eq:alternate}\\
\operatorname{First}_\theta\cap L_\theta
 &=\operatorname{First}_\theta\cap\operatorname{Hits}_\theta.
                                                        \label{eq:first}
\end{align}
Their set is effectively Presburger-definable.
\end{theorem}
\begin{proof}
Both conditions are necessary by flipping and initial colour. For sufficiency,
induct in ordinary time. A first visit has the right initial colour by
\eqref{eq:first}. Otherwise a nearest preceding visit exists: a given time has
only finitely many predecessors, even for an infinite trace. Its letter is
opposite by \eqref{eq:alternate}, so the current physical colour is correct.
The requested movement is therefore the actual movement. All integer times
are covered. Projection of violations and complementation in $D$ give the
Presburger parameter set.
\end{proof}

This quantified form of the established same-cell alternation criterion
\cite{gajardo} handles all parameter assignments together. Its advantage is that it avoids assuming an arbitrary definable
cardinality has a definable parity.

Without a supplied field, call a trace \emph{feasible} when it satisfies
the same-cell alternation condition~\eqref{eq:alternate}; it is then
realizable from some initial colouring. For a finite feasible trace, the smallest supporting
initial black set is
\begin{equation}\label{eq:root}
 B_{\min}(\theta)=V_\theta(\operatorname{First}_\theta\cap L_\theta).
\end{equation}
For an infinite trace this set may be infinite. Existence on an unspecified
colouring must not be mistaken for existence on a finite support.

\begin{theorem}[Complete endpoint]\label{thm:output}
For a finite trace, define
$\operatorname{Last}_\theta=\mathcal T_\theta\setminus\dom C_\theta$.
On its realization domain the final black field is
\begin{equation}\label{eq:out}
 B_{\rm out}(\theta)=
 \bigl(B(\theta)\setminus V_\theta(\mathcal T_\theta)\bigr)
 \cup V_\theta(\operatorname{Last}_\theta\cap R_\theta).
\end{equation}
It is an effective finite-fibre Presburger family when $B$ is one.
\end{theorem}
\begin{proof}
Unvisited initial cells retain their colour. A last right turn leaves black;
a last left turn leaves white. These cases are disjoint and exhaustive. Only
finitely many updates occur at each parameter assignment.
\end{proof}
Equation~\eqref{eq:out} retains every unread initial cell.

\section{Unbounded powers and effective resource spectra}\label{sec:power}
Consider a finite bounded-power schema
\begin{equation}\label{eq:schema}
 v_0u_1^{k_1}v_1\cdots u_s^{k_s}v_s,\qquad k_i\ge0,
\end{equation}
with fixed words and heading-restoring $u_i$. Fixed pieces may rotate heading.
For a phase $r$ in the $i$th powered block $u_i^{k_i}$, an occurrence at copy
index $j$ has time and position of the form
\begin{equation}\label{eq:affine}
 t=b_r+c_r k+\ell_r j,\qquad q=a_r+A_r k+d_r j,
 \qquad 0\le j<k_i.
\end{equation}
Here $k=(k_1,\ldots,k_s)$ is the exponent vector,
$c_r\in\mathbb Z^{1\times s}$ and $A_r\in\mathbb Z^{2\times s}$;
the remaining coefficients are fixed integer scalars or vectors of the
indicated dimensions. There are no products of two unbounded
variables. Zero powers and affine reparameterizations are allowed. A final
infinite power simply replaces its copy bound by $j\ge0$.

Theorem~\ref{thm:semantics} thus extends fixed-word replay to arbitrarily long
parameterized operations, and Theorem~\ref{thm:output} preserves the complete
field. A variable rotating power can first be split by its exponent modulo
its rotation order; our implementation instead requires this preprocessing
explicitly. Finite-fibre initial Presburger families may include finite unions
or XORs of arithmetic segments with fixed strides and variable finite lengths.

Defect end-times in a schema are Presburger: intersect shifted letter-time
sets for each allowed factor and subtract their union from eligible times.
This treats cross-block windows and zero powers exactly.

The following standard counting fact combines unary counting elimination
\cite[Theorem~5.4]{counting} with Presburger cell decomposition
\cite[Theorem~1]{cluckers}. We give the elementary effective argument needed here.
\begin{lemma}[Unary finite-fibre counting]\label{lem:count}
For an effectively Presburger set $F\subseteq\Z^d\times\Z$ with finite fibres,
$\theta\mapsto |F_\theta|$ is effectively rational-affine on the cells of a
finite Presburger partition.
\end{lemma}
\begin{proof}
Apply Presburger quantifier elimination \cite{haase}. Split parameter residues
and the counted variable modulo a common multiple of the constant moduli.
Writing $t=r+Qz$, each remaining atom either constrains the parameters alone
or supplies a floor/ceiling affine bound on $z$. Use a disjoint Boolean
partition, then partition parameters by the ordering of the finitely many bounds
and their constant-denominator residues. Every nonempty counted cell is an
integer interval with two finite rational-affine endpoints. Summing disjoint
interval lengths proves the claim. Finiteness rules out unbounded intervals
on the applicable parameter cell.
\end{proof}

\begin{theorem}[Initial-black and defect spectrum]\label{thm:spectrum}
On the feasible domain of \eqref{eq:schema}, the minimum initial-black cost
$c(\theta)=|\operatorname{First}_\theta\cap L_\theta|$ and the contained-defect
count $d(\theta)$ have effective semilinear graphs and are piecewise
rational-affine.
\end{theorem}
\begin{proof}
Apply Lemma~\ref{lem:count} to the two unary sets of times. Different first-read
times refer to different cells, so \eqref{eq:root} identifies the first count
with initial-black cost.
\end{proof}

This does not assert affine counting for arbitrary planar Presburger fibres;
for instance an $n$ by $n$ square has $n^2$ cells. The counted objects here are
one-dimensional time sets. The runtime implements exact time sets and fixed
integer thresholds via
\[
 |F_\theta|>b\iff\exists t_0<\cdots<t_b\quad\bigwedge_i t_i\in F_\theta.
\]
The counting lemma gives a constructive partition algorithm; the exact
integer-threshold tests above suffice for the supplied-schema queries.

For any supplied finite schema collection and fixed mass bound, unboundedness
of its defect spectrum is consequently decidable. If unbounded, a semilinear
component gives an affine family of finite witnesses. Such a parameter family
is not a fixed infinite counterexample.

\section{A finite cover of all prefixes reaching a fixed defect count}\label{sec:cover}
The ant-specific bridge is the existing ten-factor macro theorem
\cite{project}. We state its certificate content to make the dependency explicit.
The vertices are feasible 149-letter words whose every contained ten-factor
belongs to $\mathcal L_{10}$ (is legal); an edge appends one letter
and discards the first, provided the resulting 150-letter word is feasible
and every ten-factor is legal. The graph has 6,501 vertices and 6,605 edges. Its cyclic components have 104 (the standard word $W$), 24 (the word $T$) and 366 vertices.
The fixed words are printed in Appendix~\ref{app:words}. The 366-state
component maps to a graph with two loops: a translating word
$A$ of length 28 and a heading-reversing word $B$ of length 102. The established
physical obstruction is
\begin{equation}\label{eq:fold}
 BA^nBA^mB\text{ is infeasible for every }n,m\ge0.
\end{equation}
It is an all-parameter same-cell conflict, not an observation at small lengths.

Counting at most two completed long loops and a third incomplete loop yields
a finite covering lift with 7,595 vertices and 7,763 edges. Its only cyclic
components are simple cycles of lengths $24,28,28,28,104$. The acyclic
condensation supplies bounds of five cycle stages, 956 acyclic updates, and
212 for the sum of visited cycle lengths. Our independent reconstruction
regenerates every feasible word, checks every quotient edge, rebuilds the lift
and recovers these bounds. The exact arithmetic kernel separately verifies
\eqref{eq:fold} over the entire quadrant.

The graph construction proves the following precise coverage statement.
\begin{proposition}[Finite legal-word cover, from~\cite{project}]\label{prop:cover}
Every physically feasible finite word with no contained defect has a
representation $v_0u_1^{k_1}v_1\cdots u_r^{k_r}v_r$ with $r\le5$,
$k_i\ge0$, and $\sum_i|v_i|\le1317$. Each $u_i$ is a cyclic phase of
$W,A$ or $T$. These words restore heading. From a north-facing pose the
displayed representatives have displacements $(-2,-2),(0,-4)$ and $(0,-2)$,
respectively; cyclic phases have quarter-turn rotations of those vectors.
\end{proposition}
A finite defect-free word uses at most 149 initial letters before this graph
cover applies. Every cyclic visit is a power followed by a partial cycle and
possible exit. Therefore it has a schema with at most five powers and at most
$149+956+212=1317$ nonpowered letters. The number of ordinary updates is not
bounded: all exponents remain free.

Remove the last letter of each of the first $N$ defect windows from the
prefix ending at the $N$th defect. The at most $N+1$ remaining contiguous
pieces contain no defect: any bad window wholly in a piece would have had
its last letter removed. Apply Proposition~\ref{prop:cover} to each piece
and restore the $N$ removed letters as literal separators. Using the relaxed
constant $2048\ge1317$ gives a finite schema family with
\begin{equation}\label{eq:budget}
 s\le5(N+1),\qquad\sum_i|v_i|\le2048(N+1)+N.
\end{equation}
All power words are cyclic phases of $W,A,T$. This overapproximation may be
redundant. Exact arithmetic semantics and exact defect count are subsequently
imposed, so no nonphysical word is accepted.

\begin{proof}[Proof of Theorem~\ref{thm:main}]
For each schema in \eqref{eq:budget}, introduce its exponent variables, make
the initial field exactly the distinct coordinate list $z$, and conjoin the
realization predicate, exactly $N$ contained defects, and a defect at the last
letter. Existentially project exponents and take the finite union.
Every genuine $N$th-event prefix is covered; every satisfying assignment is a
genuine prefix by Theorem~\ref{thm:semantics}. Retain the affine time and pose
to obtain their graph. Formula~\eqref{eq:out} gives the black field; different
schemas representing the same prefix agree by determinism.
\end{proof}

The finite union supplies a total mathematical decision procedure; all
remaining parameters are eliminated by exact integer arithmetic.

\section{Six cells and arbitrarily many defects}\label{sec:six}
The resource-spectrum dichotomy has a physical fixed-mass branch. A single
six-cell shape, moved farther from the initial head, produces arbitrarily many
nonstandard windows before a certified standard future. Thus the minimum
initial mass for any finite defect target has a uniform upper bound.

Let $F$ denote one complete ant update and put
$\operatorname{trans}_v(B,p,h)=(B+v,p+v,h)$.
An intrinsic $J_{10}$ event occurs precisely when the window formed by the
current turn and the actual preceding nine turns lies outside $\mathcal L_{10}$.
The preceding turns are obtained by reversing $F$. Let $v_h$ be the unit vector
$(0,1),(1,0),(0,-1),( -1,0)$ for heading $h=\mathrm N,\mathrm E,\mathrm S,\mathrm W$,
respectively. To undo an update, move from
$p$ to $p-v_h$, undo that cell's flip, and undo the turn determined by its
restored colour. This defines the history also for windows reaching before
time zero. A forward contained defect uses only turns at nonnegative times.

\begin{theorem}[Six-cell family]\label{thm:six}
Start at the origin facing north, with complete initial black set
\begin{gather*}
 B_k=B_*+k(-2,-2),\qquad k\in\mathbb N,\\
 B_*=\{(-87,-55),(-87,-52),(-84,-55),(-82,-61),(-81,-55),(-80,-60)\}.
\end{gather*}
The entire future contains exactly $6702+47k$ forward ten-factor defects.
Its first permanent standard-P104 entrance is $18743+220k$.
The first read of an initially black cell occurs at $13318+104k$.
The same future has $6706+47k$ intrinsic $J_{10}$ events, the last at
$18751+220k$. In particular $\mu(N)\le6$ for every $N\ge1$.
\end{theorem}

The shape has $\ell^\infty$ diameter nine. Its distance from the head varies
with $k$, and its recovery after the first initial-black impact lasts
$5425+116k$ updates. We prove the parameter dependence by an exact complete-field
recurrence. All set sums in the proof below are symmetric differences, denoted
$\triangle$, unless written as translations.

\subsection{Canonical data and the incoming field}
Put $d=(2,-2)$ and use south as the canonical heading. The finite sets are
\begin{align*}
P={}&\{(0,-2),(0,-1),(0,0),(1,-1),(2,-1),(2,1),(3,-2),\\
 &\hspace{1em}(3,1),(4,-2),(4,0),(5,-1),(5,0),(6,0)\},\\
D={}&\{(0,-3),(0,0),(1,-4),(1,-3),(1,-2),(2,-1),\\
 &\hspace{1em}(2,0),(3,0),(6,-1),(7,-1),(7,0),(8,-1)\},\\
M={}&\{(-2,-8),(1,-8),(1,-5),(1,-2),(6,-1),(7,-3)\}.
\end{align*}
The literal word data $W_\star,U,V,Z,O$ have lengths $104,1210,116,4406,104$,
respectively, and are supplied in \texttt{verification/six\_cell\_feedback/witness.json}.
Here $W_\star$ is the left cyclic rotation by 57 letters of the word $W$
in Appendix~\ref{app:words}, $O$ is a cyclic phase of $W_\star$, and $V$ is the
finite return word. The blank trajectory, launched on the all-white plane from the origin facing
north, has word $C$ through 9977 updates, pose
$p_*=(-15,10),\mathrm W$, and 715 black cells. With
$\rho(x,y)=(y,-x)$, its complete field, translated by $-p_*$ and rotated by $\rho^{-1}$
to the canonical south heading, is $P\triangle E$, where $|E|=702$.
This defines $E$ by finite replay.

The exact finite and ray checks establish
\[
 F^{104}(P,(0,0),\mathrm S)=((P+d)\triangle D,d,\mathrm S),
\]
and permanent $W_\star$ futures for both $P$ and $P\triangle E$; $E$ avoids all
future reads. Consequently the complete field after $N$ standard periods,
normalized by subtracting $Nd$, is
\[
 P\triangle\mathop{\triangle}_{j=1}^{N}(D-jd)\triangle(E-Nd).
\]
For $N=32+k$, the initial markers are $p_*+N\rho(d)+\rho(M)=B_k$.
Integer-ray tests show that none belongs to the blank prefix footprint, to any
of the preceding $N$ standard periods, or to the initial inverse-nine history.
For example the standard-period condition reduces to excluding
$w_i=m+jd$ for each phase read $w_i$, $m\in M$ and integer $j\ge1$.
First-difference coupling therefore gives the physical prefix $CW_\star^N$, ending
at the complete collision field
\begin{equation}\label{eq:six-Y}
 Y_N=P\triangle M\triangle\mathop{\triangle}_{j=1}^{N}(D-jd)
          \triangle(E-Nd).
\end{equation}

\subsection{A full-field recurrence with protected archives}
Replay $U$ at $Y_{32}$. It ends at $a=(-1,-1)$, heading south. Let $Z_U$
be its odd write set, and define
\[
 K=P\triangle M\triangle Z_U,\qquad
 S_N=K\triangle\mathop{\triangle}_{j=1}^{N}(D-jd)\triangle(E-Nd).
\]
The read footprints of $U$ and of $V$ from $a$ avoid both
$\bigcup_{j\ge33}(D-jd)$ and $\bigcup_{j\ge32}(E-jd)$.
Thus $U$ executes identically for every $N\ge32$, and $V$ restores the
heading and displaces by $-d$. Let $R$ be the absolute board cells toggled an odd number of times by $V$
launched from $a$, in the same canonical coordinates as $K$ and $D$. Put
\[
 Q=K\triangle(K+d)\triangle D\triangle(R+d),\qquad |Q|=140.
\]
Expanding the preceding field formula proves, for every $N\ge33$,
\begin{equation}\label{eq:six-recurrence}
 \operatorname{trans}_{d}F^{116}(S_N,a,\mathrm S)
       =(S_{N-1}\triangle Q,a,\mathrm S).
\end{equation}
Every surviving deposit and the old blank debris appear on both sides.

The finite footprints of $V$ and of the recovery word $Z$ from $S_{32}$,
and every ray of its infinite $O$ suffix, avoid
$\mathcal A=\bigcup_{r\ge0}(Q+rd)$.
For a suffix phase $v_i$ and $q\in Q$, the last assertion reduces to
\[
 v_i+je=q+rd,\qquad j,r\ge0,\qquad e=(-2,-2).
\]
Since $e$ and $d$ are independent, determinants give the unique candidate
solution, which exact integrality and sign tests reject. The finite-point
versions use the same integer-ray equations.
Let $A_k=\mathop{\triangle}_{r=0}^{k-1}(Q+rd)$. Induction in
\eqref{eq:six-recurrence}, using the avoidance of $\mathcal A$, gives the
actual state after $k$ return words from $S_{32+k}$ as
\[
 \operatorname{trans}_{-kd}(S_{32}\triangle A_k,a,\mathrm S).
\]
The recovery word and its complete periodic future therefore execute with
all archives still present. The full infinite turn word is exactly
\begin{equation}\label{eq:six-word}
 CW_\star^{32+k}UV^kZO^\omega.
\end{equation}

\subsection{Exact counts and first entrance}
The chosen suffix boundary in \eqref{eq:six-word} is
$T(k)=9977+104(32+k)+1210+116k+4406=18921+220k$.
At $k=0$, comparison with the backward extension of $O$ puts the last
mismatch immediately before time 18743. It is inside the fixed word $Z$,
so its location shifts by $220k$. Primitivity of $W_\star$ identifies any permanent
standard suffix with the same phase. This proves the stated first entrance.

A suffix register of nine turns gives the following complete forward counts:
\begin{center}
\begin{tabular}{lr}\toprule
Block with its actual incoming suffix & Defects\\\midrule
$C$ &4537\\
First $W_\star$ after $C$ &5\\
Further $W_\star$ copies &0\\
$U$ after $W_\star$ &483\\
First $V$ after $U$ &50\\
Each further $V$ &47\\
$Z$ directly after $U$ ($k=0$)&1677\\
$Z$ after $V$ ($k\ge1$)&1674\\
$O^\omega$ after $Z$&0\\\bottomrule
\end{tabular}
\end{center}
Both parameter cases sum to $6702+47k$. The true inverse-nine history is
\texttt{LLLLRLLLL}; four additional bad windows end at times zero through
eight. Hence the intrinsic count is $6706+47k$. Its last bad window ends
eight updates after first permanent entrance, because it still includes
preceding history. The first marker read has index 13 within $U$, giving
$9977+104(32+k)+13=13318+104k$.

For any target $N\ge1$, choose
$k=\max(0,\lceil(N-6702)/47\rceil)$ and take a prefix through its $N$th
defect. This proves $\mu(N)\le6$ and realizes the unbounded resource ray
$(c,d_{10})=(6,6702)+k(0,47)$.

\paragraph{Verification of the finite premises.}
The checker reconstructs all finite words and fields, proves every required
ray avoidance, evaluates all factor seams, and checks earliest entrance.
For an infinite phase ray $q_i+ne$, the read colour is initial membership
XOR earlier-write parity. Each old cell gives at most one hit, and each
self-overlap changes parity at one integer threshold. Zero, each threshold,
each hit and its successor therefore cover every epoch, including the final
infinite interval. This is the full-drift chronological criterion used
throughout the article. The universal parameter in
\eqref{eq:six-word} follows from the recurrence and guards above.
An independent C++ implementation checks complete fields and all-future guards;
20 finite parameter controls perform 835,360 updates in each implementation.
Run \texttt{python3 defect\_synthesis/reproduce.py --six-cell-only}.


\section{Supply budgets for a finite schema collection}\label{sec:supply}
A supplied finite collection fixes its word descriptions while their exponents
remain free. Once its resource image is effectively semilinear, the following
dichotomy is a direct consequence of the signs of its period generators.
The ant-specific work is constructing the complete physical resource image;
the generator argument applies to any semilinear subset of $\mathbb N^2$.

\begin{theorem}[Supply dichotomy]
For every finite collection of bounded-power schemas, exactly one of the
following alternatives holds, and the alternative is effectively decidable:
\begin{enumerate}
\item Some fixed initial cost $c_0$ supports arbitrarily many contained defects.
One can extract a schema and an affine ray of exponents
$k(n)=k_0+nv$ with $c(k(n))=c_0$ and $d(k(n))=d_0+n\delta$, $\delta>0$.
\item There are computable rational $\alpha\ge0,\beta$ such that
$d(w)\le\alpha c(w)+\beta$ for every feasible instance in the collection.
\end{enumerate}
\end{theorem}
\begin{proof}
Project the semilinear resource graph to $(c,d)\in\N^2$ and compute a finite
union of linear sets $b+\sum_j\N g_j$. A generator with zero cost coordinate
and positive defect coordinate gives the first alternative. Conversely, if
cost stays bounded while defects grow, positive-cost generators have bounded
multiplicities in some component, so a positive vertical generator is necessary.
When there is none, set
\[
 \alpha=\max\bigl(\{0\}\cup\{g_{j,d}/g_{j,c}:g_{j,c}>0\}\bigr),\qquad
 \beta=\max_b(b_d-\alpha b_c).
\]
The inequality holds at each base and is preserved by every generator.
Use zero constants for the empty spectrum. To extract the first alternative's
exponents, fix its $c_0$ in the full $(k,c,d)$ graph. An unbounded component
has a generator with positive $d$ coordinate, giving the asserted affine ray.
The minimal seeds are computed by first-left reads for each parameter.
\end{proof}
The ray describes a family of finite inputs. The constants in the other
alternative belong to the specified collection. When a collection has
unbounded defects and no vertical ray, its least supply for at least $N$
defects is $\Theta(N)$: the affine budget gives the lower bound, and any
positive-defect generator gives the upper bound.

\subsection{Three independent fold and consumer parameters}
Use the inherited 28-letter word $A$, 102-letter connector $B$, and
\[P=\mathrm{RLRRRRLRRRRLLLLRLRRL}.\]
From $(-1,0)$ facing west consider
\[
 w(a,b,r)=B[9:]A^aBA^bP^r,\qquad a,b,r\ge0.
\]
Here $B[9:]$ deletes nine initial letters. The input field is synthesized
from first reads, so the three exponents are independently quantified.

\begin{theorem}[Exact fold--consumer spectrum]\label{thm:supply}
The word is feasible exactly when
\[
 a\ge1,\qquad 0\le b<a,\qquad 0\le r\le b+2.
\]
On this domain, its minimum initial-black cost and contained-defect count are
\[
 C=46+5a+5b,\qquad
 D=\begin{cases}0&r=0,\\9r-3&r\ge1.\end{cases}
\]
Its resource image is exactly
\begin{align*}
 &(51,0)+\N(5,0),\\
 &(51,6)+\N(5,0),\\
 &(51,15)+\N(10,9)+\N(5,0).
\end{align*}
Consequently $10D\le9C-309$, with sharp slope and intercept. For a complete
input of mass $m\ge51$, the largest possible count within this family is
\[
 D_{\max}(m)=15+9\left\lfloor\frac{m-51}{10}\right\rfloor.
\]
\end{theorem}
\begin{proof}
Compile every phase of the five blocks into its exact position--time relation.
Same-cell successor alternation gives a Presburger feasible set; the supplement
checks its equality with the displayed region, with no bound on $a,b,r$.
First-left times are partitioned into disjoint residue intervals. Each proposed
interval is accepted only after exact set equality with the corresponding
fibre, so its integer length counts all first-left reads. Their sum is
$46+5a+5b$. The shifted ten-letter matcher includes all junction windows and
zero powers, and its time-set count gives the stated $D$. These are exact
symbolic identities in the supplied finite certificate.

For $r=0,1$, put $b=0,a=1+s$ to obtain the first two components. For $r\ge2$,
put $r=q+2$; its least feasible $b,a$ are $q,q+1$. Increasing $a$ supplies
all extra multiples of five, yielding precisely the third component.
The inequality holds on all three components and is equality at
$a=b+1,r=b+2$ for every $b\ge0$, proving both sharpness claims.
For the mass optimum take $b=\lfloor(m-51)/10\rfloor$, $a=b+1,r=b+2$;
place any remaining black cells outside the finite read footprint.
\end{proof}
This extends the symmetric producer family used in the factor paper to all
independently sized folds and finite consumer trains. Its terminal black mass
is $67+9(a+b)+4r$, so initial supply and final mass measure different resources.


\section{Small positive witnesses and unary NP membership}\label{sec:small}
Semilinearity provides a decision procedure. The fixed geometry of the cycle
library also gives a quantitative size bound on a positive certificate.
Let $E(m,N)$ mean that some complete seed of at most $m$ black cells produces
at least $N$ contained defects, with all coordinates unrestricted. Put
\[
 s=5(N+1),\quad K=2570(N+1),\quad d=s+2m,\quad
 H=2^{24}(K+d+1),\quad R=64\bigl((d+1)!H^d+1\bigr).
\]
\begin{theorem}[Small witnesses]\label{thm:small}
If $E(m,N)$ holds, it has a witness whose initial coordinates and schema
exponents have absolute value at most $R$. Its $N$th defect occurs within
$K(R+1)$ updates. In particular,
$\log R=O((m+N+1)\log(m+N+2))$.
With $m,N$ in unary, $E(m,N)$ belongs to NP.
\end{theorem}
The radius is existential: it produces a possibly different input with the
same budgets. A specified input may put its next obstacle arbitrarily far away.

\subsection{A polynomial verifier for a compressed physical run}
The canonical cover has at most $s$ powers and $2048(N+1)+N$ literal letters.
There are at most $K$ phase types: the literal bound plus at most
$104\cdot5(N+1)$ powered phases is
$2568(N+1)+N\le2570(N+1)$. A powered type has position
$A_i(k)+jD_i$ and affine time, with $0\le j<k_i$. The nonzero strides,
after rotations, are
\[
 (\pm2,\pm2),\quad(\pm2,0),\quad(0,\pm2),\quad
 (\pm4,0),\quad(0,\pm4).
\]
For a proposed phase occurrence, its read colour is initial membership XOR
the parity of every earlier type visiting the same cell. Each nonzero-stride
type visits a given cell at most once. Solve
$A+jD=B+\ell E$ for each pair of types. Nonparallel strides give at most
one integral pair, with determinant dividing 16. Parallel strides yield a
bounded interval in one residue class of $j$, of modulus at most two.
A source point or literal type gives at most one hit.

Different blocks have their fixed chronological order. Within one powered
block, intersecting phases have $\ell=j+c$ for a constant $c$; the sign of
$|u|c+\text{old phase}-\text{current phase}$ decides precedence. Thus all
previous flips are included. Split the current index by its hit moduli,
sort interval starts and one-past-ends, and check the parity on each nonempty
interval against the requested bit. Induction on ordinary update time proves
that these checks are necessary and sufficient for execution on the complete
supplied field.

There are $O(K(K+m))$ pair and point queries and
$O(K(K+m)\log(K+m)+m^2)$ comparisons, including source distinctness.
All operate on binary integers of polynomial length. A word summary stores
length, first and last nine symbols, and contained-defect count; concatenation
checks only crossing windows and binary powering is logarithmic. In this
library all powered words are cyclically ten-legal and have length at least
24, so the count is fixed once every exponent is classified as zero or positive.

\subsection{Linear chambers through successful certificates}
Let $\theta$ list the independent exponents and the coordinates of $m'\le m$
distinct initial black cells; its dimension is at most $d$. At one successful
witness fix $\theta\bmod64=\rho$ and write $\theta=64z+\rho$. Follow the verifier
and record every comparison with its witnessed outcome, including endpoint
ordering, equalities, source distinctness, and exponent zero/positive cases.
Strict integer inequalities become non-strict inequalities by a unit shift.
For a failed equality record its strict sign.

All modular tests are now constant. Determinant denominators divide 16,
point-membership denominators divide four, and the parallel-index and parity
splits each divide by at most two; the accumulated divisors divide 32.
Floors and ceilings are affine in $z$, with only the constant term rounded.
Every integer solution to the resulting chamber follows the same comparison
path, hence preserves physical feasibility and exact defect count.

For completeness, the coefficient envelope can be followed explicitly. Set
$H_0=2^8(K+d+1)$. After the substitution, origin variable coefficients are at
most 256 and their constants at most $K+252s$. Origin differences fit
$2H_0$. Determinant numerators fit $16H_0$ and their quotients fit
$16H_0+1$. Parallel-hit endpoints fit $4H_0+16$; after parity splitting and
successors all endpoints fit $16H_0+32$. Minima and maxima select one existing
form. Endpoint comparisons and collinearity guards fit $128H_0+128$.
Within-body origin differences cancel all parameters. Thus every integer
coefficient and right-hand side is bounded by $256H_0\le H$.

\begin{proof}[Proof of Theorem~\ref{thm:small}]
The integer small-solution theorem gives a solution to every nonempty
$n$-variable chamber of coefficient height $a\ge1$ with coordinates at most
$(n+1)!a^n$ in absolute value~\cite{small}. Apply it with $n\le d,a\le H$,
then substitute $\theta=64z+\rho$. This proves the stated $R$.
The source still has distinct cells, the powers remain nonnegative, and its
word still ends at the $N$th defect. Summing at most $K$ phase types gives the
duration bound. Taking logarithms yields the bit bound.

For unary budgets guess the schema, source coordinates and exponents inside
these bounds. They use $O((m+N+1)^2\log(m+N+2))$ bits. Check the conservative
schema-size limits, library membership, exact physical execution and defect
count using the preceding algorithm. Every accepted certificate is a genuine
run, and the finite cover plus small-solution argument supplies one for every
positive instance. This proves NP membership.
\end{proof}
The chamber argument applies to the canonical library with free initial
coordinates. A supplied affine input family with larger coefficients requires
its own coefficient-height bound.

\begin{corollary}\label{cor:mu}
The minimum initial mass $\mu(N)$ needed for at least $N$ contained defects
can be computed in deterministic time $2^{O(N^2\log(N+2))}$.
\end{corollary}
\begin{proof}
The word $(\Ll\Rr\Rr\Ll)^q$ has minimum supply $2q$ and
$\max(0,4q-9)$ defects, so $\mu(N)\le2\lceil(N+9)/4\rceil=O(N)$.
Enumerate the bounded positive certificates for successive masses through
this upper bound. The preceding bit-size estimate gives the claimed time.
\end{proof}
Short schemas and small integer solutions are established tools for flat
counter systems~\cite{demri}. Here the mutable-field interval verifier supplies
the additional ant-specific execution argument.


\section{A polynomial-size quantifier-free decision backend}\label{sec:qf}
The general Presburger construction supports arbitrary finite-fibre source
predicates and complete target equality. For the modal schemas in the
finite-defect cover, a narrower syntax admits a polynomial-size existential
encoding. This provides an executable decision backend and a second route
to the small-witness conclusion.

Quantifier-free encodings of counting properties on flat runs already appear
in Decker--Pirogov~\cite{flat}. Here the additional obligation is to encode
all spatial coincidences and chronologically earlier flips in a complete
mutable field. The restricted strides below make those parity constraints
polynomial in the supplied schema size.

Let
\[
 \mathcal D=\{(\pm2,\pm2),(\pm2,0),(0,\pm2),(\pm4,0),(0,\pm4)\}.
\]
Every powered word must restore heading and have full drift in $\mathcal D$.
The source is a finite XOR of segments $a+jv$, $0\le j<L$, with
$v\in\mathcal D$, and singletons. Anchors, lengths and exponents may be
affine expressions in shared integer parameters; lengths and exponents
are constrained nonnegative. Explicit $m$-point inputs additionally require
pairwise distinct coordinates. No coordinate or repetition bound is imposed.

\begin{theorem}[Quantifier-free physical realization]
For this syntax, a polynomial-time construction produces a quantifier-free
linear-integer/Boolean circuit whose existential satisfiability is equivalent
to literal execution of the complete prescribed word. Exact contained-defect
counts, the final defect flag, endpoint time and pose, and any polynomial
number of endpoint-colour queries have polynomial-size circuits as well.
\end{theorem}
\begin{proof}
Represent each phase of a power by the injective read stream
$a_i+j d$, $0\le j<k$, with its fixed required colour. Its actual colour
is the XOR of the source and every earlier planned write at that point.
For streams in different blocks, temporal order is fixed. Within one
power of period $P$, a phase-$r$ occurrence at index $j+s$ precedes
phase $i$ at index $j$ exactly when $Ps+r<i$. This retains the intra-loop
chronology rather than grouping all occurrences by phase.

For a source or write stream $b+\ell v$, solve
$a+jd=b+\ell v$ with its finite index bounds. Nonparallel strides give at
most one hit by the determinant formula. Parallel strides give a guarded
arithmetic interval in $j$ with modulus
$|v_\alpha|/\gcd(|d_\alpha|,|v_\alpha|)$, which is always 1 or 2 for
$\mathcal D$. Splitting $j=2z+e$ leaves a guarded ordinary interval.
Floors and exact divisibility use only fixed integer divisors and therefore
have constant-size linear defining constraints.

On a target residue interval $[0,K]$, the XOR of these interval contributions
is a step function. It has the required constant colour exactly when its
value at zero is correct and every active interior boundary has even total
jump parity. Check this at each candidate interval start and one-past-end,
including all coincident boundaries. For $q$ contributions, the equality
and parity tests use $O(q^2)$ shared wire operations; no ordering of symbolic
endpoints is enumerated. With $R$ read streams and $M$ source terms the
parity stage has size $O(R(R+M)^2)$, apart from shared affine expressions.
Induction through the finite planned word proves that these colour tests
are sufficient as well as necessary for physical execution.

For lexical counts, retain the preceding at most nine turns, their length,
the defect count and the final-event flag as shared registers. For a power
$u^k$, at most $\lceil9/|u|\rceil$ copies precede suffix stabilization;
later copies contribute a fixed cyclic defect count. Constant-size guarded
transfers select the short cases and the stabilized case. The $k=0$ branch
preserves the entire incoming state. Sharing these transfers between blocks
gives a circuit linear in spelled length and block count, instead of a
Cartesian list of zero/positive choices.

Finally, an endpoint-colour query XORs initial membership with membership
in every read stream. This retains all unqueried remote cells. Named gates
and linear definitions preserve polynomial size throughout and prove the
theorem.
\end{proof}

\begin{corollary}[Unary budgets with binary side conditions]
For explicit $m$ coordinate pairs, unary $N\ge1$ and a quantifier-free
linear-integer side condition $Q$ of binary length $b$, existence of a
complete $m$-black source whose $N$th contained defect satisfies $Q$ is
in NP. Positive witnesses have polynomial bit length in $m+N+b$.
For unary $N$, the minimum initial cost $\mu(N)$ is in $\mathrm{FP}^{\mathrm{NP}}$.
\end{corollary}
\begin{proof}
The inherited complete cover supplies a schema with at most $5(N+1)$
powers and $2048(N+1)+N$ fixed letters. Its powered bodies are phases of
$W,T,A$, with strides in $\mathcal D$. Guess this polynomial-length
schema, apply the theorem to the $m$ distinct source points, and require
exactly $N$ defects with a bad last letter and condition $Q$.
Guessing Boolean branches reduces the circuit to polynomially many linear
integer constraints. The integer small-solution bound used in
Section~\ref{sec:small} supplies polynomial-bit witnesses, which can be
checked in polynomial time. This uses the complete physical cover;
satisfiability of a graph relaxation alone would not suffice.

Theorem~\ref{thm:six} gives $\mu(N)\le6$. A finite prefix can be padded
by unread black cells outside its footprint, so the NP query for exactly $m$
points decides $\mu(N)\le m$ without additional side conditions. Searching
$m\in\{0,\ldots,6\}$ uses a constant number of NP-oracle queries, each of
polynomial unary-$N$ size, and proves the stated function-class bound.
\end{proof}

The encoding applies to the stated segment-source syntax and endpoint
colour queries. The general compiler handles quantified source predicates,
zero-displacement powers and complete-field target equality. Polynomial circuit size does not
imply a polynomial deterministic solver, and the NP statement is not a
claim for binary-encoded $N$.

The reproduction archive checks the encoding by exact intersection
identities and complete finite endpoint comparisons. Its Z3 and ISL/GMP
backends are part of the computational evidence described in the
reproduction section.


\section{Finite targets and coordinate dependence}
The minimum initial mass introduced in Section~1 can equivalently be
computed from finite feasible words:
\[
 \mu(N)=\min\{c(w):w\text{ is finite feasible and }d(w)\ge N\}.
\]
Indeed every input supplies the first-left reads of each prefix, and every
feasible word has the minimal seed in~\eqref{eq:root}.
The six-cell construction in Theorem~\ref{thm:six} gives $0\le\mu(N)\le6$.
Theorem~\ref{thm:main} makes every fixed-mass query decidable, so this finite
upper bound proves that $\mu$ is a total computable nondecreasing function.
Corollary~\ref{cor:mu} supplies a quantitative computation bound for this function.

Theorem~\ref{thm:six} places a point in $E_{6,N}$ for every $N$. The chosen point varies with the target: each constructed input has finitely many defects and then enters P104. Thus the fixed-mass resource ray has a complete positive dynamical realization.

The synthesis theorem fixes the requested event count and quantifies over
all coordinates. The supply dichotomy instead fixes a finite schema collection
and allows unbounded repetitions. Keeping these two inputs explicit gives a
way to study which physical mechanisms can produce many defects at low cost.

\section{Limitations}
The bound $\mu(N)\le6$ is a uniform upper bound; minimality of six is a separate
question. The six-cell family moves its initial shape relative to the head as
$k$ grows. Each member has its own finite defect total and a permanent P104 tail.

The all-coordinate construction takes fixed numerical budgets $m,N$;
its formulas need not have uniformly bounded size when these budgets vary.
The supply dichotomy applies to a supplied finite schema collection. Its
sharp constants describe that collection, and a growing family of finite
witnesses is distinct from a single infinite trajectory. Fixed-prefix
decision procedures alone do not settle eventual convergence.

The polynomial-size quantifier-free backend assumes the heading-restoring
stride library and source syntax specified in Section~\ref{sec:qf}; the
complexity statements use unary budgets. Unary time-fibre counting is
essential: general planar Presburger fibres can have nonlinear cardinality.
The proof gives an effective universal disjunction for nonzero budgets;
the reported computations verify the finite cover and supplied schemas
rather than expand every instance of that disjunction. Exact programs and
written reductions constitute the computational proof, with language and
arithmetic-library implementations in its trust base.

\section{Reproducibility}
Run \texttt{python3 defect\_synthesis/reproduce.py}. The self-contained
supplement contains two exact relation compilers, independent finite-state
controls, the finite-cover reconstruction, the complete three-parameter
resource identity and fixed-schema query optimizations. Python 3.10+, a C++17
compiler and the system ISL/GMP and Z3 shared libraries are required. The arithmetic
backend performs exact integer set operations~\cite{isl}; the retained
computational reports use ISL 0.27.

The accompanying source package is bound by its explicit \texttt{sources.json} manifest; it includes the six-cell word data, recurrence checks and independent C++ replay.
Its \texttt{sources.json} lists the exact source and certificate hashes.
The imported factor results are Theorems A--B (attraction), Subsection 3.3
(the all-parameter fold obstruction), and Theorem F with its finite
covering lift in Subsections 4.1--4.2 of the archived article~\cite{project}.

The programs check supplied-schema identities and reconstruct the finite cover.
The universal nonzero-budget disjunction is established by the constructive
proof; its full expansion is not needed by these computations. The small-witness
theorem follows from the interval and chamber argument in Section~\ref{sec:small}.

For repeated integer threshold queries, the optimized kernel represents a
unary time set by its strict-order relation. Its $m$th relation power witnesses
$m+1$ distinct times, so binary powering reduces a fresh threshold query to
$\lfloor\log_2m\rfloor+\operatorname{popcount}(m)-1$ compositions for $m\ge1$.
A shared decision graph also evaluates the 84 allowed ten-factors with 183
mask intersections instead of 840. Exact set-equality controls compare the
optimized queries with their baseline semantics. These improvements preserve
the mathematical schema cover and the complete endpoint field.

\appendix
\section{Fixed turn words}\label{app:words}
The following representatives use the origin and north-facing initial heading,
with N,E,S,W step vectors $(0,1),(1,0),(0,-1),(-1,0)$.
Line breaks carry no meaning. A cyclic phase is a rotation of the letters;
its displacement is a quarter-turn rotation of the displayed representative's
vector when the initial heading is reset to north.

\noindent $W$ (104 letters):
\begin{quote}\ttfamily\small
RRLLLLRLRRRRLRRRRLLLLRLRRRRLRLLRRLLLLRRLLRRRRLLRRLRL\\
LRLLRLRLLRRRRLLRLLRRRRLLRRRRLLRLRRRRLRLLLLRRRRLRRLRR
\end{quote}

\noindent $T$ (24 letters):
\begin{quote}\ttfamily\small
RRRRLLRLLRRRRLLRRRRLLRLL
\end{quote}

\noindent $A$ (28 letters):
\begin{quote}\ttfamily\small
LLRRRRLRRLRRRRLLLLRLRRRRLRLL
\end{quote}

\noindent $B$ (102 letters):
\begin{quote}\ttfamily\small
RRLLLLRRLLRRRRLLRRLRLLRLLRLRLLRRRRLLRLRRRRLRLLLLRRRR\\
LRRLRRRRLLLLRLRRRRLRLLLLRRRRLRRLRRRRLLLLRLRRRRLRLL
\end{quote}

\clearpage
\begingroup\small
\begin{thebibliography}{10}\setlength{\itemsep}{1pt}
\bibitem{gajardo} A.~Gajardo. A symbolic projection of Langton's Ant.
\emph{Discrete Mathematics \& Theoretical Computer Science}, DMCS'03, 2003.
\href{https://doi.org/10.46298/dmtcs.2312}{doi:10.46298/dmtcs.2312}.
\bibitem{haase} C.~Haase, S.~N.~Krishna, K.~Madnani, O.~S.~Mishra and
G.~Zetzsche. An efficient quantifier elimination procedure for Presburger
arithmetic. \emph{ICALP 2024}, LIPIcs 297, 142:1--142:17.
\href{https://doi.org/10.4230/LIPIcs.ICALP.2024.142}{doi:10.4230/LIPIcs.ICALP.2024.142}.
\bibitem{flat} N.~Decker and A.~Pirogov. Flat model checking for counting LTL
using quantifier-free Presburger arithmetic. \emph{VMCAI 2019},
LNCS 11388, 513--534, 2019.
\href{https://arxiv.org/abs/1901.05692}{arXiv:1901.05692}.
\bibitem{isl} S.~Verdoolaege et al. \emph{Integer Set Library manual},
ISL 0.27, \texttt{doc/user.pod} in the
\href{https://libisl.sourceforge.io/isl-0.27.tar.xz}{ISL 0.27 source distribution}.
\bibitem{project} H.~Ke. \emph{Local Factor Rigidity and Mass--Defect Bounds
for Langton's Ant}. Version 0.1.0, 2026.
\href{https://doi.org/10.5281/zenodo.23665478}{doi:10.5281/zenodo.23665478}.
\bibitem{small} R.~Bottesch, M.~W.~Haslbeck, A.~Reynaud and R.~Thiemann.
Verifying a solver for linear mixed integer arithmetic in Isabelle/HOL.
\emph{NASA Formal Methods}, LNCS 12229, 233--250, 2020.
\href{https://doi.org/10.1007/978-3-030-55754-6_14}{doi:10.1007/978-3-030-55754-6\_14}.
\bibitem{demri} S.~Demri, A.~K.~Dhar and A.~Sangnier.
Taming past LTL and flat counter systems.
\emph{Information and Computation} 242, 306--339, 2015.
\href{https://doi.org/10.1016/j.ic.2015.03.007}{doi:10.1016/j.ic.2015.03.007}.
\bibitem{ginsburg} S.~Ginsburg and E.~H.~Spanier.
Semigroups, Presburger formulas, and languages.
\emph{Pacific Journal of Mathematics} 16(2), 285--296, 1966.
\href{https://doi.org/10.2140/pjm.1966.16.285}{doi:10.2140/pjm.1966.16.285}.
\bibitem{counting} N.~Schweikardt. Arithmetic, first-order logic, and counting
quantifiers. \emph{ACM Transactions on Computational Logic} 6(3), 634--671, 2005.
\href{https://doi.org/10.1145/1071596.1071602}{doi:10.1145/1071596.1071602}.
Author version: \href{https://arxiv.org/abs/cs/0211022}{arXiv:cs/0211022}.
\bibitem{cluckers} R.~Cluckers. Presburger sets and $p$-minimal fields.
\emph{Journal of Symbolic Logic} 68(1), 153--162, 2003.
\href{https://doi.org/10.2178/jsl/1045861509}{doi:10.2178/jsl/1045861509}.
\end{thebibliography}
\endgroup
\end{document}
